% ============================================================
% Capitulo 6 -- Discussao e Conclusao
% Tamanho-alvo: 4-6 paginas | Sem figuras
% Ordem de escrita: QUARTO (depois de Resultados e Aplicacao)
% ============================================================

\chapter{Discuss\~ao e Conclus\~ao}
\label{cap:conclusao}

% ============================================================
\section{S\'intese dos achados por pergunta de pesquisa}
\label{sec:sintese}

% TESE: cada RQ tem resposta clara, ancorada em secao do Cap 4.
%
% CONTEUDO (~1.5 pagina) -- estrutura por RQ:
%
% RQ1 (sinal acima do textual?):
% - FNN: texto da raiz satura sinal (LogReg=0.86, §\ref{sec:baseline_textual}).
%   GNN adiciona < 2pp.
% - UPFD-GossipCop: topologia agrega valor (+6pp via SAGE,
%   §\ref{sec:explainer}); refinado em §\ref{sec:estratificacao} para
%   +13pp em depth>=3 e +21pp em depth>=5.
% - UPFD-PolitiFact: GNN nao agrega valor significativo.
%
% RQ2 (sob quais condicoes?):
% - Dominios com Cohen's d >= 0.5 exibem sinal explorivel
%   (§\ref{sec:porque}). |d| < 0.2 nao.
%
% RQ3 (so topologia e' viavel?):
% - Sim, condicional ao d acima. Parte b: invariancia confirmada
%   PT vs EN; parcialmente refutada DE vs EN, causa nao isolada
%   (§\ref{sec:multilingual}).

% ============================================================
\section{Achado negativo \'e achado}
\label{sec:negativo}

% TESE: PolitiFact nao e' falha de implementacao -- e' observacao
%       cientifica sobre aplicabilidade do paradigma.
%
% CONTEUDO (~1 pagina):
% - Cohen's d desprezivel em PolitiFact (§\ref{sec:porque}) mostra que
%   fake/real sao estruturalmente indistinguiveis.
% - Reportar com a mesma extensao dos positivos = exigencia da tese central.

% ============================================================
\section{Limita\c{c}\~oes}
\label{sec:limitacoes}

% TESE: 7 limitacoes especificas, todas reconhecidas e nao mascaradas.
%
% CONTEUDO (~1 pagina) -- lista numerada:
%
% (a) FakeNewsNet com cascatas estrela (script 00) por falta de
%     timestamps precisos. NOTA: UPFD-GossipCop *oficial* nao e
%     estrela plana -- 30% tem depth >= 3 (§\ref{sec:estrat_cascatas}).
%     Esta limitacao vale para o FNN que construimos, nao para UPFD.
% (b) Regra dual 0.8/0.2 derivada de UM experimento (UPFD-GossipCop).
% (c) Bluesky sem labels = so demonstracao (Cap 5).
% (d) Reprodutibilidade nao bit-exact (CUDA/cuDNN).
% (e) Sem comparacao numerica direta com BiGCN/GCNFN -- so via valores
%     publicados.
% (f) Bluesky e snapshot 2024-2025, nao amostra probabilistica --
%     generalizacao a janelas futuras nao garantida.
% (g) CPU-only (Python 3.12.10, torch 2.10+cpu) -- hiperparametros
%     podem ser suboptimos em GPU.

% ============================================================
\section{Trabalhos futuros}
\label{sec:futuros}

% TESE: 4 direcoes especificas e factiveis.
%
% CONTEUDO (~1 pagina) -- 4 itens:
%
% 1. Pseudo-labels no Bluesky usando LogReg-BERT-FNN como oraculo,
%    retreino especializado.
% 2. Calibracao de probabilidades (Platt scaling / isotonica) com
%    ground-truth de algum subset.
% 3. Expansao para datasets PolitiFact-like recentes (Twitter pos-2023,
%    Mastodon).
% 4. Replicacao direta de BiGCN \cite{bian2020bigcn} / GCNFN para
%    benchmark numerico contra a literatura.

% ============================================================
\section{Considera\c{c}\~oes finais}
\label{sec:final}

% TESE: contribuicao = diagnostico do paradigma, nao nova arquitetura.
%       Defesa lado a lado de positivo e negativo e' o metodo central.
%
% CONTEUDO (~0.5 pagina, 1 paragrafo):
% Sem grandiloquencia. Texto ancora:
% "Este trabalho nao propos uma nova arquitetura GNN. Propos um
%  diagnostico -- sob quais condicoes a topologia da propagacao detecta
%  fake news, e quando nao detecta. A defesa de um resultado negativo
%  (PolitiFact) lado a lado de um resultado positivo (GossipCop),
%  ancorada em estatistica de tamanho de efeito, e a contribuicao
%  metodologica que esperamos transcender o escopo do dataset especifico."
